\section{Methods}
\label{sec:methods}

This section summarizes (i) the audit metrics used to quantify and localize cross-lens mismatch, (ii) the common-staging control, and (iii) the cosmological inference setup used to estimate $\mnu$ posteriors in the presence of lens swaps.

\subsection{Directional cross-lens audits and localization}
\label{sec:audit-metrics}

Our primary audit statistic is the directional held-out mismatch $\dchi_{B|A}$ defined in \cref{eq:dchi_def}. Intuitively, $\dchi_{B|A}$ asks: if we package the macro object using lens $A$, how poorly does that packaged solution predict lens $B$ relative to $B$'s own internal best fit? Because route mismatch is directional, we report both $\dchi_{B|A}$ and $\dchi_{A|B}$.

When the likelihood supports blockwise decomposition (e.g.\ by spectrum and multipole band), we additionally localize mismatch by writing $\dchi_{B|A}$ as a sum over blocks. In our implementation, blocks correspond to a spectrum type (TT/TE/EE) and an $\ell$-bin group. The result is a heatmap-style diagnostic that turns a global mismatch into a concrete localization: which spectrum and which multipole ranges dominate the held-out failure.

\paragraph{Blockwise accounting.}
When a likelihood provides additive blockwise contributions (e.g.\ by spectrum and multipole groups), we write
\begin{equation}
  \log \mathcal{L}_B(\theta) \;=\; \sum_{g\in\mathcal{G}} \log \mathcal{L}_{B,g}(\theta),
  \label{eq:block_decomp}
\end{equation}
and define the blockwise held-out mismatch contributions
\begin{equation}
  \dchi_{B|A,g} \;\equiv\; -2\left[\log \mathcal{L}_{B,g}(\hat\theta_A) - \log \mathcal{L}_{B,g}(\hat\theta_B)\right].
  \label{eq:block_dchi}
\end{equation}
Under \cref{eq:block_decomp} these satisfy the sum rule $\dchi_{B|A}=\sum_{g\in\mathcal{G}}\dchi_{B|A,g}$ (and similarly for $A|B$), so the localization heatmaps should be read as an accounting decomposition of the same directional audit statistic rather than an unrelated diagnostic.
We emphasize that $\dchi$ is reported in the native $-2\Delta\log\mathcal{L}$ units of the packaged likelihood objects and is used here as a comparative diagnostic (across directions and staging controls), not as a calibrated hypothesis-test p-value; null calibration would require additional held-out or posterior-predictive constructions beyond the scope of this short note.

\subsection{Profiled audits: separating shared parameters from nuisance packaging}
\label{sec:profiled-audit}

A key subtlety is that lenses can differ in nuisance parameterization and calibration degrees of freedom. To more closely align the audit with SBT's ``completion'' semantics, we also implement a \emph{profiled} variant of the cross-lens audit. For direction $B|A$, we:
\begin{enumerate}[leftmargin=*, itemsep=2pt]
  \item hold fixed a small set of \emph{shared} parameters (the intersection of parameters we choose to treat as common across lenses), and
  \item re-optimize a tractable subset of $B$'s nuisance parameters (a calibration/bandpower-mode subset) to maximize $\log \mathcal{L}_B$ under those shared parameters.
\end{enumerate}
We then compute a profiled $\dchi_{B|A}$ by comparing $B$'s likelihood at the profiled $A$-shared point against $B$'s own profiled best fit. Operationally, this separates (i) mismatch that can be absorbed by lens-specific nuisance packaging from (ii) mismatch that persists even after giving the test lens freedom to retune its nuisance degrees of freedom.

In practice, full profiling over all nuisance parameters can be computationally expensive. We therefore use a runtime-capped profiling set consisting of a calibration/beta subset (referred to internally as \texttt{TT\_cal\_beta}), reduced to a tractable number of parameters for deterministic optimization.
\paragraph{Formal definition.}
Let parameters be partitioned as $\theta=(\phi,\eta_B)$ where $\phi$ denotes the shared parameters we hold fixed in the audit direction $B|A$ and $\eta_B$ denotes nuisance parameters specific to lens $B$. Define the profiled log-likelihood
\begin{equation}
  \log \widetilde{\mathcal{L}}_B(\phi) \;\equiv\; \max_{\eta_B}\, \log \mathcal{L}_B(\phi,\eta_B),
  \label{eq:profile_loglike}
\end{equation}
and the profiled directional mismatch
\begin{equation}
  \widetilde{\dchi}_{B|A} \;\equiv\; -2\left[\log \widetilde{\mathcal{L}}_B(\hat\phi_A) - \log \widetilde{\mathcal{L}}_B(\hat\phi_B)\right],
  \label{eq:profile_dchi}
\end{equation}
where $\hat\phi_L$ is the shared-parameter component of the best-fit point under lens $L$.
In practice we approximate \cref{eq:profile_loglike} by profiling only a tractable subset of $B$'s nuisance parameters under a runtime cap; the intent is not to ``profile away'' all mismatch, but to reduce sensitivity to purely lens-specific packaging degrees of freedom.
In our implementation, the profiled set comprises the following calibration/beta parameters:
\begin{quote}\small\raggedright
\texttt{Tcal\_ext150},
\texttt{Tcal\_rel90},
\texttt{Tcal\_rel220},
\texttt{Ecal\_ext150},
\texttt{Ecal\_rel90},
\texttt{Ecal\_rel220},
\texttt{TT\_CIBClustering\_Alpha},
\texttt{TT\_CIB\_150x150},
\texttt{TT\_CIB\_150x220},
\texttt{TT\_CIB\_220x220},
\texttt{TT\_GalCirrus\_Alpha},
\texttt{TT\_GalCirrus\_Amp}.
\end{quote}
\noindent For the reverse direction, the capped set is:
\begin{quote}\small\raggedright
\texttt{TT\_tSZ\_CIB\_Corr\_Amp},
\texttt{TT\_tSZ\_CIB\_corr},
\texttt{Ecal150},
\texttt{Ecal220},
\texttt{Ecal90},
\texttt{TT\_CIBClustering\_Alpha},
\texttt{TT\_CIBClustering\_Amp},
\texttt{TT\_CIBClustering\_Beta},
\texttt{TT\_GalCirrus\_Alpha},
\texttt{TT\_GalCirrus\_Amp},
\texttt{TT\_GalCirrus\_Beta},
\texttt{TT\_Poisson\_150x150}
\end{quote}
\noindent (runtime-capped to 12 parameters; see run metadata).

\subsection{Common-staging control}
\label{sec:common-staging}

Because the two SPT lenses differ in multipole support, a basic staging control is to restrict the broader-support lens to the narrower-support range and repeat the same audits. We implement a \emph{common-staging} control in which the D1 lens is restricted (as much as feasible) to the SPT2018 multipole support (TT: $\ell\ge 750$; TE/EE: $\ell\ge 300$; all with $\ell\le 3000$). The restriction is implemented at likelihood-instantiation time via selection masks recorded in the run metadata. If mismatch persists under common-staging and remains localized inside the overlap region, this points away from a pure support/cut explanation.
This common-staging control matches \emph{multipole support} but does not fully equalize the compressed summaries: the two releases retain their own bandpower binning/window functions and internal covariance conventions, so this test should be interpreted as a support/cut control rather than an identity-of-data-vector comparison.

\subsection{$\mnu$ inference pipeline}
\label{sec:mnu-inference}

We estimate $\mnu$ posteriors using Cobaya \cite{cobaya} with a standard Boltzmann backend (CLASS \cite{class} and/or CAMB \cite{camb}, depending on the configuration). The canonical SPT-only+DESI results use \textbf{CAMB} while the Planck-subset baseline+DESI results use \textbf{CLASS}; see run metadata for exact configs. We sample the standard six $\Lambda$CDM parameters $\{\omega_b h^2, \omega_c h^2, H_0, \ln(10^{10}A_s), n_s, \tau\}$ together with $\mnu$ under a physical gate $\mnu\ge 0$. Following the neutrino modeling description in \cite{arxiv2601_16277}, our $\Lambda$CDM+$\mnu$ completion assumes three degenerate massive neutrino states when $\mnu$ is allowed to vary.

\paragraph{Implementation caveat (CLASS parameterization).}
CLASS expects massive neutrinos to be specified via $N_{\mathrm{ncdm}}$ and per-species masses $m_{\mathrm{ncdm}}$. To keep $\mnu$ explicitly present in the recorded chains and summaries while providing CLASS the required per-species masses, we sample an internal parameter (denoted \texttt{mnu\_sample}) and define a derived parameter $\mnu$ together with $m_{\mathrm{ncdm}}=(\mnu/3,\mnu/3,\mnu/3)$. This is a parameterization workaround; it preserves the reported $\mnu$ posterior but should be kept in mind when comparing config syntax to other pipelines.

\paragraph{Priors and stability checks.}
The lower bound $\mnu\ge 0$ is explicit in \cite{arxiv2601_16277}. The upper bound of the prior is not stated in the paper source we audited; for reproducibility we therefore use a broad flat prior with upper bound $5\,\mathrm{eV}$ and report this choice explicitly. Since the resulting posteriors and 95\% upper bounds we report are all $\ll 1,\mathrm{eV}$, this upper bound is not active in practice; the relevant gate for boundary behavior is the physical constraint $\mnu\ge 0$. For each production run we record basic chain diagnostics (split-$\hat R$ \cite{gelman_rubin1992,vehtari2021_rhat} and a crude ESS estimate \cite{vehtari2021_rhat} for $\mnu$) and a half-chain stability proxy for the 95\% upper bound. We also perform a mild stability check by repeating (or partially re-running) one production configuration with a distinct random seed and verifying that the 95\% upper bound does not change materially.
Posterior summaries (quantiles, KDE-based 1D densities, and simple convergence heuristics) are computed using GetDist \cite{getdist}.

\subsection{Reproducibility note}
\label{sec:reproducibility}

All results in this note are tied to canonical run bundles recorded in \path{docs/findings/freeze_checklist.md} and summarized in \path{docs/findings/canonical_results.json}.
All figures and tables appearing in the paper are curated from, or auto-generated against, canonical run bundles and a canonical results summary, so the numerical claims in the text are backed by a deterministic artifact pipeline rather than manual transcription.
The full analysis code, experiment runners, and manuscript build workflow are available at \url{https://github.com/ioannist/six-birds-neutrino}.
The figures curated under \path{paper/figures/} and the tables under \path{paper/tables/} are generated from these canonical artifacts. For convenience, the primary run bundles cited later include:
\begin{itemize}[leftmargin=*, itemsep=2pt]
  \item SPT-only + DESI $\mnu$ shift comparison: \path{runs/20260212_080120_648405_mnu_shift_spt_only_desi}.
  \item Planck-subset baseline + DESI $\mnu$ shift comparison: \path{runs/20260212_111357_576743_mnu_shift_cmb_plancksubset_desi}.
  \item Profiled baseline vs common-staging audit comparison: \path{runs/20260211_181258_958327_spt_profiled_common_staging_compare}.
\end{itemize}
A one-command rebuild of key artifacts is provided by \texttt{make freeze\_check}, and the neutrino inference workflows are exposed via \texttt{make mnu\_smoke} and \texttt{make mnu\_full} when required external assets are present.
